\documentclass[12pt]{article}
\usepackage[margin=1in]{geometry}
\usepackage{amsmath}
\usepackage{amssymb}
\usepackage{enumitem}
\usepackage{hyperref}

\begin{document}

\begin{center}
\textbf{SOSC 13200 (W26): Social Science Inquiry II --- Practice In-Class Final}
\end{center}

\noindent\textbf{Time:} 80 minutes \\
\textbf{Notes:} Closed notes. No R / no computers. \\
\textbf{Allowed:} Pencil/pen; non-programmable calculator. \\
\textbf{Instructions:} Answer all questions. Show your work for calculations; partial credit is available. If you use notation, define it.

\vspace{0.5em}
\noindent\textbf{Provided formulas (you may use these)}
\begin{itemize}
  \item Independence: $\Pr[A\cap B]=\Pr[A]\Pr[B]$
  \item 95\% CI (normal approx): estimate $\pm 1.96\times SE$
  \item $t$-statistic (generic): $t=\dfrac{\text{estimate} - \text{null}}{SE}$
\end{itemize}

\noindent\textbf{Power (two-arm design, equal $n$ per arm; normal approximation)}
\begin{itemize}
  \item Solve for $n$:
  \[
  n \approx 2\left(\frac{(z_{1-\alpha/2}+z_{\text{power}})\sigma}{\text{MDE}}\right)^2
  \]
  \item Notation: $z_p$ is the $p$th quantile of $N(0,1)$ (so $\Pr[Z\le z_p]=p$ for $Z\sim N(0,1)$).
  \item Helpful quantiles (given): $z_{0.975}=1.96$, $z_{0.80}=0.84$
  \item If a pre-treatment covariate explains fraction $R^2$ of outcome variance, residual sd is approximately:
  \[
  \sigma_{\text{resid}} \approx \sigma\sqrt{1-R^2}.
  \]
\end{itemize}

\clearpage
\noindent\textbf{Question 1 (6 points): Hypothesis tests, p-values, and confidence intervals}

\noindent\textbf{Setup:} A city runs a randomized ``get-out-the-vote'' experiment. Some registered voters receive a door-to-door
canvassing visit ($W=1$) and others do not ($W=0$). Let $Y$ be an indicator for whether a person votes in the next election.
Suppose your estimate of the difference in mean turnout between treated and control voters is:
\begin{itemize}
  \item $\widehat{ATE}=0.05$ (a 5 percentage-point increase)
  \item Standard error: $SE(\widehat{ATE})=0.03$
\end{itemize}

\begin{enumerate}[label=1\alph*)]
  \item (2 pts) In words, state what the null hypothesis $H_0: ATE=0$ means in this setting.
  Then compute the $t$-statistic for testing it.
  \item (2 pts) Compute an approximate 95\% confidence interval for the ATE.
  \item (2 pts) Briefly interpret the confidence interval and what it implies about statistical significance at the 5\% level.
\end{enumerate}

\vspace{1.0em}
\noindent\textbf{Question 2 (6 points): Regression and covariate adjustment}

\noindent\textbf{Setup:} You study the effect of a job-training program on monthly earnings.
Program participation $W$ is randomized ($W=1$ offered training, $W=0$ not offered). You also have a \textbf{baseline} measure
of monthly earnings, $base$, collected before randomization.

\vspace{0.5em}
\noindent You estimate two regressions (standard errors in parentheses):

\vspace{0.25em}
\noindent\textbf{Model 1 (unadjusted):} $Y_i=\alpha + \tau W_i + \varepsilon_i$ \\
$\hat\tau=120\ (80)$

\vspace{0.5em}
\noindent\textbf{Model 2 (adjusted):} $Y_i=\alpha + \tau W_i + \gamma\, base_i + \varepsilon_i$ \\
$\hat\tau=120\ (50)$ \\
$\hat\gamma=0.70\ (0.06)$

\vspace{0.75em}
\begin{enumerate}[label=2\alph*)]
  \item (3 pts) Interpret $\hat\tau$ from Model 1 in words.

  \vspace{0.5em}
  \item (3 pts) In a randomized experiment, how is Model 1 related to a difference-in-means estimator?
  Then explain briefly why adding $base$ can reduce the standard error on $\hat\tau$.
\end{enumerate}

\vspace{1.0em}
\noindent\textbf{Question 3 (4 points): Randomization inference (permutation test)}

\noindent\textbf{Setup:} A researcher studies discrimination in email responses from landlords.
They randomize the perceived race of a prospective tenant by varying the name on an email inquiry ($W=1$ vs $W=0$),
and record whether the landlord replies ($Y=1$ reply, $0$ otherwise).
Suppose the observed difference-in-means estimate is $\widehat{ATE}_{obs}=-0.07$.
They run a randomization inference test by shuffling the treatment label $W$ (without replacement) $B=2000$ times, recomputing a placebo ATE each time.

\noindent They find that in \textbf{150} of the 2000 permutations, $|\widehat{ATE}^{(b)}|\ge |\widehat{ATE}_{obs}|$.

\begin{enumerate}[label=3\alph*)]
  \item (2 pts) Compute the (two-sided) randomization-inference p-value.
  \item (2 pts) Briefly interpret this p-value in words.
\end{enumerate}

\vspace{1.0em}
\noindent\textbf{Question 4 (4 points): Power, effect sizes, and sample size}

\noindent\textbf{Setup:} You are designing a two-arm follow-up experiment with equal sample size per arm,
outcome sd $\sigma=1.0$, two-sided $\alpha=0.05$, and target power $=0.80$.

\begin{enumerate}[label=4\alph*)]
  \item (2 pts) Use the provided power formula to compute the approximate sample size \textbf{per arm} needed to detect $\text{MDE}=0.20$.
  \textbf{Show your steps} by filling in the intermediate quantities (use the helpful quantiles given above):
  \begin{itemize}
    \item $z_{1-\alpha/2}=z_{0.975}=\underline{\hspace{2cm}}$
    \item $z_{\text{power}}=z_{0.80}=\underline{\hspace{2cm}}$
    \item $z_{1-\alpha/2}+z_{\text{power}}=\underline{\hspace{2cm}}$
    \item $\left(\dfrac{(z_{1-\alpha/2}+z_{\text{power}})\sigma}{\text{MDE}}\right)=\underline{\hspace{2cm}}$
    \item $n\approx 2\left(\dfrac{(z_{1-\alpha/2}+z_{\text{power}})\sigma}{\text{MDE}}\right)^2=\underline{\hspace{2cm}}$
  \end{itemize}

  \vspace{0.75em}
  \item (2 pts) Now suppose you will adjust for a strong pre-treatment covariate, and you expect it to explain $R^2=0.50$ of the outcome variance.
  \begin{enumerate}[label=\arabic*.]
    \item (1 pt) Approximate the new residual standard deviation $\sigma_{\text{resid}}$.
    \item (1 pt) Plug $\sigma_{\text{resid}}$ into the same power formula to get the new required $n$ per arm (approximately).
    In one sentence: does the required sample size go up, go down, or stay the same?
  \end{enumerate}
\end{enumerate}

\vspace{1.0em}
\noindent\textbf{Question 5 (4 points): Observational regression and causal claims}

\noindent\textbf{Setup:} You have observational data on individuals. Let $X$ be weekly hours of social media use and
$Y$ be a mental-health index (higher values mean better mental health).
You estimate the bivariate regression:
\[
Y_i=\alpha + \beta X_i + \varepsilon_i,
\]
and you find $\hat\beta=-0.30$.

\begin{enumerate}[label=5\alph*)]
  \item (2 pts) Give one plausible confounder (an omitted variable) that could make $\hat\beta$ differ from the causal effect
  of social media use on mental health. State whether you think the bias is likely \textbf{upward} or \textbf{downward}, and explain briefly.

  \vspace{0.5em}
  \item (2 pts) Suppose your sample is extremely large, and you estimate $\hat\beta=-0.30$ with a tiny standard error
  (so the p-value is essentially 0). Explain why this does \textbf{not} by itself justify a causal interpretation
  of $\hat\beta$, and name one additional design feature or kind of evidence that would make a causal claim more credible.
\end{enumerate}

\vspace{1.0em}
\noindent\textbf{Question 6 (6 points): Inference, measurement, and preregistration}

\begin{enumerate}[label=6\alph*)]
  \item (4 pts) In 3--5 bullet points, sketch a \textbf{pre-analysis plan} for a randomized experiment that has
  multiple treatment arms plus a control (e.g., different versions of a reminder message).
  Your plan must include:
  \begin{enumerate}[label=\arabic*.]
    \item one \textbf{primary estimand/contrast} (state it in words), and
    \item the \textbf{null hypothesis} for that primary contrast (state it in words), and
    \item two additional things you would commit to in advance to reduce researcher degrees of freedom
    (e.g., handling missing data, covariates, subgroup analyses, multiple comparisons, outliers).
  \end{enumerate}

  \vspace{0.75em}
  \item (2 pts) A researcher measures 12 outcomes and reports only the outcomes with $p<0.05$.
  \begin{enumerate}[label=\arabic*.]
    \item In one sentence: what credibility problem does this create?
    \item Suppose (for simplicity) the 12 tests are independent and \textbf{all null hypotheses are true}.
    What is the probability of seeing \textbf{at least one} statistically significant result at $\alpha=0.05$?
  \end{enumerate}
\end{enumerate}

\end{document}

